\title{
\includegraphics[width=0.3\textwidth]{example-image}~\\[0.5cm]
নিবন্ধ সংকলন}
\author{\footnotesize{সম্পাদনায়} \\ 
সম্পাদক\\ 
সহ-সম্পাদক}

\date{\small ২০২১}

\begin{titlepage}
\includepdf[pages=1, fitpaper]{cover}
\thispagestyle{empty}
\predate{\centering \footnotesize{প্রকাশনায়}\\ \large{\href{<url>}{প্রকাশক}} \\}
\postdate{\vfill\hfill \footnotesize{\\অলংকরণে\\ মুবতাসিম ফুয়াদ\\  \uccoff{\fontfamily{lmr}\selectfont ©} \href{https://github.com/rafisics/ebook-template}{আদীব হাসানের টেমপ্লেট}\\}\hfill}     % https://tex.stackexchange.com/a/199356/114006
\maketitle
\end{titlepage}

% Adding dedication and preface                     https://tex.stackexchange.com/a/51539/114006

\pagestyle{empty}
\vspace*{\stretch{1}}
\begin{flushright}
\itshape
আমাদের \\ প্রত্যেক পাঠক এবং শুভাকাঙ্ক্ষীর \\ প্রতি উৎসর্গীকৃত
\end{flushright}
\vspace{\stretch{3}}

\cleardoublepage

\newpage
\pagestyle{plain}
\pagenumbering{roman}

\chapter{সম্পাদকীয়}

		\chapquote{\qtext{One glance at a book and you hear the voice of another person, perhaps someone dead for 1000 years. To read is to voyage through time.}}{কার্ল সেগান}

		\firstword{T}{he} is the preface/abstract of the book.
		\blindtext \\

		\section*{কৃতজ্ঞতা}
		\firstword{T}{he} is the acknowledgement of the book.
		\blindtext \\

		\begin{flushright}
		\itshape
		সম্পাদক, \href{<url>}{নিবন্ধ সংকলন} \\ জুন ২০২১
		\end{flushright}

\cleardoublepage

\setcounter{tocdepth}{0}
\tableofcontents

\newpage
\pagestyle{plain}
\pagenumbering{arabic}